\bumppoints{28}

\begin{multicols}{2}
\parbox{\linewidth}{\question
Boreal forests in Canada contain large evergreen trees like spruce trees and pine trees. If a boreal forest continues to get about 140 cm of precipitation annually but mean annual temperature increases from 0°C to 5°C over the next 100 years, the ecosystem will

\medskip

\begin{choices}
\choice remain a boreal forest. 
\choice become more like a temperate rain forest. 
\choice become more like a tundra. 
\correctchoice become more like a temperate deciduous forest. 
\choice become more like a sclerophyllous woodland. 
\end{choices}}

\columnbreak
	\includegraphics[width=\linewidth]{exam4_climograph2}
\end{multicols}

\question
Considering your knowledge of the dead zone and the causes of microscopic algae blooms, the Gulf of Mexico ecosystem is \emph{most likely} to be limited by which of the following nutrients?

\begin{choices}
\choice hydrogen
\correctchoice nitrogen
\choice carbon
\choice \ce{CO2}
\choice fertilizer
\end{choices}

\question
Which of these species is most likely a keystone species in Alaska?

\begin{choices}
\choice Kelp.  
\choice Sea urchins.  
\choice Humans.  
\correctchoice Sea otters. 
\choice Killer whales. 
\end{choices}

\question
On average, the amount of secondary consumer biomass is about \underline{\hspace{2cm}} less than the amount of primary producer biomass. 
\begin{choices}
\choice 10$\times$ %dontshuffle 
\choice 20$\times$ %dontshuffle 
\choice 50$\times$ %dontshuffle 
\choice 80$\times$ %dontshuffle 
\correctchoice 100$\times$%dontshuffle 
\end{choices}


\question
The energy stored by producers can be passed through an ecosystem along a(n) \rule{1in}{0.4pt}, a series of interactions in which organisms transfer energy by eating and being eaten.

\begin{choices}
\choice nutrient cycle 
\choice ecosystem 
\correctchoice food web 
\choice trophic cascade 
\choice predator-prey chain  
\choice energy cycle 
\end{choices}

\question
A tertiary consumer is one trophic level above secondary consumer. The tertiary consumer has high abundance. In a trophic cascade, which of the following is probably \emph{not} true.
\begin{choices}
\choice The secondary producer will have high abundance. 
\correctchoice The primary producer will have high abundance. 
\choice The decomposer level will have low abundance. 
\choice The secondary consumer will have low abundance. 
\choice The primary consumer will have high abundance. 
\end{choices}



\question
The Mediterranean coast and coastal California have similar types of ecosystems. This is because

\begin{choices}
\choice they are both wine-growing regions.
\correctchoice they have similar climates.
\choice they have similar types of communities.
\choice they have similar types of organisms.
\choice they are at similar latitudes.
\end{choices}


\question
Based on the hypothetical food web below, which letter represents an organism that is most likely a decomposer?  The arrows point to organisms receiving energy.

\begin{multicols}{2} 
	\begin{choices}
		\correctchoice A %dontshuffle %EndChoice
		\choice B %dontshuffle %EndChoice
		\choice C %dontshuffle %EndChoice
		\choice D %dontshuffle %EndChoice
		\choice E %dontshuffle %EndChoice
	\end{choices}
\columnbreak
	\includegraphics{exam4_foodweb}
\end{multicols}
\vspace{-1\baselineskip}

\question
A fish ate shrimp. The shrimp ate algae (which are photosynthetic).  The algae is a \rule{1in}{0.4pt}. The fish is a \rule{1in}{0.4pt}.

\begin{choices}
\choice primary consumer; primary producer. 
\choice secondary consumer; primary consumer. 
\choice secondary consumer; primary producer. 
\choice primary producer; primary consumer. 
\choice primary consumer; secondary consumer. 
\correctchoice primary producer; secondary consumer. 
\end{choices}


\question
Consider the food chain grass $\to$ grasshopper $\to$ mouse $\to$ snake $\to$ hawk. How much of the chemical energy fixed by photosynthesis of the grass (100\%) is available to the hawk?

\begin{choices}
\correctchoice 0.01\%
\choice 0.1\% %dontshuffle 
\choice 1\% %dontshuffle 
\choice 10\% %dontshuffle 
\choice 60\% %dontshuffle 
\end{choices}


\question
Intermediate successional species modify the environment, which allows late successional species to colonize. This is known as

\begin{choices}
\choice primary succession 
\choice secondary succession 
\choice tolerance. 
\choice inhibition 
\choice tertiary succession 
\choice disturbance 
\correctchoice facilitation. 
\end{choices}

\question
The dead zone in the Gulf of Mexico and other oceans around the globe occurs when

\begin{choices}
\choice eutrophication causes an absence of nutrients in the zone.
\correctchoice an excess of nutrients causes an algal bloom, followed by death of the algae. Decomposition of the algae uses up the oxygen in the water, so other organisms can no longer live there.
\choice commercial fishers over-harvest shrimp, crab, fishes and other sea animals.
\choice global climate change causes the water to become too warm to sustain life.
\choice an excess of nutrients causes an algal bloom. The algae consume all of the oxygen in the water for photosynthesis, so other organisms can no longer live there.
\end{choices}

\question
After transient orca began to prey on sea otters, what happened to the kelp forest?

\begin{choices}
\choice The kelp abundance decreased. Sea otters and kelp have a symbiotic relationship. Without the sea otters, the symbiotic relationship was broken and so the kelp abundance decreased.
\choice The kelp abundance increased. Otters dislodge kelp with their playful antics.  Because there were fewer otters, less kelp was dislodged so the kelp abundance increased.
\choice The kelp abundance remained about the same. Only the predator-prey interactions between Orca and otters were important, so nothing happened to the kelp.
\choice The kelp abundance remained about the same. The kelp are not part of this community, so the kelp could not be affected.
\correctchoice The kelp abundance decreased. Sea otters eat urchins and urchins eat kelp.  Because there were fewer otters to eat urchins, the number of urchins increased.  The urchins therefore ate more kelp so the kelp abundance decreased.
\end{choices}

\newpage

\question
In ecosystems, why is ``cycling'' used to describe the transfer of chemical elements and nutrients, while ``flow'' is used for the transfer of energy?

\begin{choices}
\choice Energy flows from one ecosystem to other ecosystems, but chemical elements constantly cycle within an ecosystem.
\choice Both chemical elements and energy flow within an ecosystem but can cycle to other ecosystems.
\correctchoice Chemical elements are recycled within ecosystems, but energy flows through and out of ecosystems.
\choice Both chemical elements and energy are recycled but can flow to other ecosystems.
\choice Chemical elements are cycled into ecosystems from other ecosystems, but energy constantly flows within the ecosystem.
\end{choices}
